Let $(\Sigma,g)$ be a closed connected oriented hyperbolic (constant curvature $-1$) surface of genus $\textsl{g} \geq 2$. Let $L \to \Sigma$ be a Hermitian line bundle equipped with a unitary connection $\nabla$, and denote by $F_\nabla = -i B \vol \in C^\infty(\Sigma, \Lambda^2 T^*\Sigma)$ the curvature $2$-form of $\nabla$, where $B \in C^\infty(\Sigma)$, and $\vol$ is the Riemannian volume. We~call~$B$ the \emph{magnetic field}. The \emph{magnetic Laplacian} is defined as
\begin{equation}
\label{equation:delta}
\Delta_L := \tfrac{1}{2} \nabla^* \nabla : C^\infty(\Sigma,L) \to C^\infty(\Sigma,L).
\end{equation}
More generally, taking $L^{\otimes k}$ for $k \in \Z_{\geq 0}$ and the induced connection $\nabla^{\otimes k}$, one can form similarly to \eqref{equation:delta} a Laplacian $\Delta_k$ acting on $C^\infty(\Sigma,L^{\otimes k})$. {\color{black}Note that the curvature on $L^{\otimes k}$ is $-ikB\vol$.} The operator $k^{-2}\Delta_k$ is a twisted semiclassical (pseudo)differential operator, with semiclassical parameter $h := k^{-1} > 0$.

Throughout this article, we~will further assume that $B$ is \emph{constant}. Note that, by~Gauss-Bonnet, this implies that $2B(\textsl{g}-1) \in \Z$. The purpose of this paper is to study the semiclassical limits of the Laplace eigenstates
\begin{equation}
\label{equation:eigenstates}
k^{-2}\Delta_k u_k = (E+\eps_k)u_k, \qquad u_k \in C^\infty(\Sigma,L^{\otimes
k}), \; \|u_k\|_{L^2(\Sigma,L^{\otimes k})} = 1,
\end{equation}
in the regime $k \to +\infty$, where $E \geq 0$ and $\eps_k \to_{k \to +\infty} 0$. {\color{black}This means that we consider the regime where the magnetic strength $kB$ (that is, the curvature of $L^{\otimes k}$) tends to infinity.} We say that $u_k$ converges to the semiclassical defect measure $\mu$ (defined on~$T^*\Sigma$) if for all $a \in C^\infty(T^*\Sigma)$,
\begin{equation}
\label{equation:convergence}
\langle\Op_{k}(a)u_{k},u_{k}\rangle_{L^2} \underset{k \to \infty}\to \int_{T^*\Sigma} a(x,\xi) \dd\mu(x,\xi),
\end{equation}
where $\Op$ is a certain semiclassical magnetic quantization on $\Sigma$ (see Section~\ref{section:microlocal} for a definition). We~denote by $u_k \rightharpoonup_{k \to \infty} \mu$ this convergence. By the normalization assumption $\|u_k\|_{L^2(\Sigma,L^{\otimes k})} = 1$, $\mu$ is a probability measure; in addition, $\mu$ must be supported on the compact energy shell $\{p=E\} \subset T^*\Sigma$ (where $p(x,\xi) := \tfrac{1}{2}|\xi|^2$). The limit \eqref{equation:convergence} may not exist as $k \to +\infty$; however, it~always exists along a subsequence $(k_n)_{n \geq 0}$ and we will often drop the subsequence notation. We~shall see that three different semiclassical regimes appear according to the value of $E$, affecting the possible behaviour of $\mu$.

\subsection{Main results}

Let $\omega_0$ be the Liouville symplectic $2$-form on $T^*\Sigma$. Set
\[
\Omega := \omega_0 + i \pi^* F_\nabla,
\]
where $\pi : T^*\Sigma \to \Sigma$ is the projection. This is the Liouville $2$-form with a magnetic correction; it is still a symplectic $2$-form on $T^*\Sigma$. The (semiclassical) principal symbol of $k^{-2}\Delta_k$ is $p(x,\xi) = \tfrac{1}{2}|\xi|^2_g$. The Hamiltonian vector field $H_p^\Omega$ of $p$ computed with respect to $\Omega$ is defined through the relation
\[
dp(\csbullet) = \Omega(\csbullet, H^{\Omega}_p).
\]
The Hamiltonian flow $(\Phi_t)_{t \in \R}$ generated by $H_p^{\Omega}$ is called the \emph{magnetic flow}. As any autonomous Hamiltonian flow, $(\Phi_t)_{t \in \R}$ preserves the (compact) energy layers $\{p=E\}$ and a natural smooth probability measure $\mu_{\mathrm{Liouv}}$ on $\{p=E\}$ called the Liouville measure (see Section~\ref{ssection:geo} for these definitions).

Let $S\Sigma \to \Sigma$ be the unit tangent bundle. Denote by $(\varphi_t)_{t \in \R}$ the (Anosov) geodesic flow, $(R_t)_{t \in \R}$ the $2\pi$-periodic rotation in the circle fibers of $S\Sigma$, and $(h_t)_{t \in \R}$ the stable horocyclic flow. Define
\[
E_c := \tfrac{1}{2}B^2, \quad T_E = (B^2-2E)^{-1/2} \text{ if } E < E_c, \quad T_E := (2E-B^2)^{-1/2} \text{ if } E > E_c.
\]
The energy $E_c$ is called the \emph{critical energy}. The following fact is
well-known \cite{Arnold-61,Ginzburg-96} and will be reproved quickly in
Section~\ref{ssection:magnetic-flow}. For $E > 0$, the magnetic flow $(\Phi_t)_{t \in \R}$ on $\{p=E\}
\subset T^*\Sigma$ is conjugate either to:
\begin{enumerate}

\item the reparametrized rotation flow $(R_{t/T_E})_{t \in \R}$ of period
$2\pi \cdot T_E$ if $E < E_c$ (elliptic case),

\item the horocyclic flow $(h_t)_{t \in \R}$ if $E=E_c$ (parabolic case),

\item the reparametrized geodesic flow $(\varphi_{t/T_E})_{t \in \R}$ if $E>E_c$ (hyperbolic case).
\end{enumerate}
Notice that for $E < E_c$ the space of orbits at energy $E$ is diffeomorphic to $\Si$ itself, so~the $(\Phi_t)_{t \in \R}$ invariant probability measures of $\{ p = E\}$
identify with the probability measures of $\Si$. This holds as well for
$E=0$ because the flow is stationary on $\{ p =0 \}$.

We will prove that this transition in the dynamics of the magnetic flow at the critical energy $E=E_c$ translates at the quantum level into the following result on semiclassical defect measures:

\begin{theorem}[The three classical/quantum regimes]
\label{theorem:main}
The following holds under the above assumptions on $(\Sigma,g)$ and $(L,\nabla)$:
\begin{enumerate}

\item\label{theorem:maini} \emph{Low energies regime.} If $0 \leq E < E_c$, then for any $(\Phi_t)_{t \in \R}$ invariant probability measure $\mu$ on $\{p=E\}$, there exists a sequence $(u_{k})_{k \geq 0}$ satisfying \eqref{equation:eigenstates} such that $u_{k} \rightharpoonup_{k \to \infty} \mu$;

\item\label{theorem:mainii} \emph{Critical energy regime.} For any
sequence $(u_k)_{k \geq 0}$ satisfying \eqref{equation:eigenstates} with $E=E_c$, $u_{k} \rightharpoonup_{k \to \infty} \mu_{\mathrm{Liouv}}$

\item\label{theorem:mainiii} \emph{High energies regime.} If $E_c < a < b$, consider for
all $k \geq 0$, the set of eigenstates $(u_{k,j})_{k,j \in A}$ such that
$k^{-2} \Delta_k u_{k,j} = \lambda_{k,j} u_{k,j}$ with $\lambda_{k,j} \in
[a,b]$ and $A \subset \Z_{\geq 0}^2$. Then there exists a density one subset
$A_\star \subset A$ such that for all sequences $(u_{k_n,j_n})_{n \geq 0}$
with $(k_n,j_n) \in A_\star$ and $k_n \to \infty$, $u_{k_n,j_n} \rightharpoonup_{n \to \infty} \mu_{\mathrm{Liouv}}$.
\end{enumerate}
\end{theorem}
